% Drawing 152/01-7 -- folder 152, batch 01, archive page 1 (PDF page 2).
% Transcription: en haut à droite du feuillet, quatre cercles : un cercle dont partent sept rayons depuis le centre, avec trois arcs du bord épaissis ; un petit cercle d'où rayonnent huit traits ; un disque hachuré portant cinq sommets marqués, de quatre desquels part un faisceau de deux ou trois traits ; plus bas, un cercle plus grand avec trois départs de traits, dont un double en haut à droite.
% Kind: four circles with radii, marked vertices and bundles of rays. Method 2 (schematic redraw from the facsimile). Status: ok.
% Notes: Four figures. (1) Left: a circle with four strokes from a point near its centre (one up to the rim, one to the rim at the upper right, ending at a heavy dot, one to the lower left and one to the lower right, both running a little past the rim) and three heavy arcs on the rim (upper left down the left side, bottom, upper right down the right side), with a heavy dot at the top. (2) Top right: a small circle with a double heavy tick at its top and another at its upper right; rays: one to the upper left, one to the left, one to the upper right, one to the right, three from its foot; faint pencil strokes: a cross inside, a stroke at the lower left and one at the lower right crossing the rim, a long faint one at the upper right; a short stroke above at the left. (3) Below it: a closed outline (not a true circle: heavy from the left vertex over the top to the upper right vertex) with five heavy vertices (top, upper right, left, lower right, bottom) and five oblique pencil strokes hatching it; rays: two from the upper right (and a short tick beside them), two from the left, three from the bottom, one long one from the lower right with a double stroke above it, and a double tick under the lower right vertex. (4) Bottom: an oval with two strokes leaving one point of its rim at the upper right (and a short third stroke there), one long stroke to the left and one stroke downwards from its foot. Left out: the long faint pencil lines that cross figure (1) and run over the sheet, the pencil star with two blackened wedges at its left, a few isolated ink dots round figure (3), and everything else on the sheet (the other figures of this trial sheet are other drawings). Overlaid on the 400 dpi colour facsimile on 2026-11-12 (Opus 5.5) and corrected: in (1) three of the seven radii of the earlier redraw were the long pencil lines passing through the centre, they are removed (the « sept rayons » of the transcription count those pencil lines; four strokes are inked from the centre), the circle and its heavy arcs are refitted (1 mm larger); in (2) the strokes are redrawn one by one where they stand (the earlier redraw had a faint pencil line as an ink ray and had missed the upper right double tick); in (3) the earlier redraw had three rays at the upper right and three at the left where the page has two and two, and a ray of the bottom in the wrong place; the outline, the five hatching strokes and every ray are redrawn from the page; in (4) the oval and its strokes refitted (1 to 2 mm). Redrawn 2026-11-12 (issue 128): to check again.
% Coordinates: 1 unit = 1 mm = 1/100 of the crop width (bbox in page fractions: [0.745, 0.085, 0.985, 0.27]).
% Redrawn by Claude (claude-fable-5-2), issue #128, 2026-11-12; not checked by a human; not compiled (no TeX in the session).
% Corrected by Claude (claude-opus-5-5) 2026-11-12 (issue #128): overlaid on the 400 dpi colour facsimile and corrected from it. Not compiled. Status: ok.
% Rights: reproduces part of the work -- behind RENDER_TRANSCRIPTS, never published by itself.
\documentclass[tikz,border=3pt]{standalone}
\usepackage{amsmath,amssymb}
\usetikzlibrary{patterns,arrows.meta}
\tikzset{pencil/.style={black!45}, struck/.style={black!35}, mpt/.style={circle,draw,inner sep=1.3pt}, pt/.style={circle,fill,inner sep=0.9pt}, opt/.style={circle,draw,fill=white,inner sep=1.1pt}}
\begin{document}
\begin{tikzpicture}[x=1mm,y=1mm,line cap=round,line join=round,font=\scriptsize]
\draw (16.6,-37.8) circle[radius=12.5];
\draw (16.2,-37.9) -- (17.3,-25.5);
\draw (16.2,-37.9) -- (21.2,-26.2);
\draw (16.2,-37.9) -- (3.3,-44.2);
\draw (16.2,-37.9) -- (27.4,-45.5);
\draw[very thick] (12.3,-26.1) arc[start angle=110, end angle=199, x radius=12.5, y radius=12.5];
\draw[very thick] (6.2,-44.8) arc[start angle=214, end angle=310, x radius=12.5, y radius=12.5];
\draw[very thick] (28.8,-40.4) arc[start angle=-12, end angle=57, x radius=12.5, y radius=12.5];
\node[pt] at (19.8,-25.3) {};
\draw (62.4,-11.7) ellipse[x radius=5.2, y radius=5.6];
\draw (51.3,-3.3) -- (57.6,-9.8);
\draw (51.5,-11.2) -- (57.3,-10.4);
\draw (67.2,-9.8) -- (73.4,-6.4);
\draw (67.6,-10.4) -- (77.8,-10.3);
\draw (61.3,-17.2) -- (57.8,-22);
\draw (61.3,-17.2) -- (61.6,-21.8);
\draw (61.3,-17.2) -- (65.3,-22.2);
\draw[thick] (61.1,-3.3) -- (60.9,-8.2);
\draw[thick] (62.2,-3.5) -- (62.2,-7.5);
\draw[thick] (67.3,-5.3) -- (66.4,-8.1);
\draw[thick] (68.4,-6.1) -- (67.2,-8.5);
\draw (58.7,-4.4) -- (59.1,-5.9);
\draw[pencil] (64.5,-7.6) -- (69.6,-3);
\draw[pencil] (54.8,-19.6) -- (57.4,-16.3);
\draw[pencil] (64.4,-14.6) -- (67.1,-17.4);
\draw[pencil] (61,-9.3) -- (63.4,-13.3);
\draw[pencil] (63.6,-10) -- (61.3,-11.9);
\draw[thick] plot[smooth, tension=0.55] coordinates {(55.8,-39.9) (55.5,-36) (56.5,-34) (58.9,-32) (61.5,-30.4) (64.8,-29.7) (68.5,-30.4) (72.1,-32) (75.9,-33.9)};
\draw plot[smooth, tension=0.55] coordinates {(75.9,-33.9) (75.6,-37) (75.2,-40.5) (74.75,-44.4) (73.5,-46) (70,-48.2) (64.7,-49.4) (59.5,-48.3) (56.6,-46) (55.5,-42.5) (55.8,-39.9)};
\draw[pencil] (61.8,-31.2) -- (56.8,-38.5);
\draw[pencil] (67,-31.5) -- (58.7,-44.6);
\draw[pencil] (69.5,-32.5) -- (61.5,-45.3);
\draw[pencil] (72,-35) -- (65.3,-46);
\draw[pencil] (73,-39) -- (68.7,-48);
\node[pt] at (64.8,-29.7) {};
\node[pt] at (75.9,-33.9) {};
\node[pt] at (55.8,-39.9) {};
\node[pt] at (74.75,-44.4) {};
\node[pt] at (64.7,-49.4) {};
\draw (75.9,-33.9) -- (88.75,-23);
\draw (75.9,-33.9) -- (92,-30.3);
\draw (77.8,-32.7) -- (78.8,-32.2);
\draw (55.8,-39.9) -- (46.75,-35.2);
\draw (55.8,-39.9) -- (45.3,-42);
\draw (64.7,-49.4) -- (60.6,-57.9);
\draw (64.7,-49.4) -- (65.3,-60.5);
\draw (64.7,-49.4) -- (71.4,-56.3);
\draw (74.75,-44.4) -- (86.7,-51);
\draw (76.1,-43.9) -- (80.9,-41.25);
\draw (76.4,-42.9) -- (80.6,-40.5);
\draw (74.1,-46.6) -- (74.75,-50.9);
\draw (75.1,-46.75) -- (75.75,-50.4);
\draw (55.85,-91.1) ellipse[x radius=14.45, y radius=13.7];
\node[pt] at (63.8,-79.6) {};
\draw (63.8,-79.6) -- (69.9,-62.4);
\draw (63.8,-79.6) -- (76.4,-63.5);
\draw (63.8,-79.6) -- (66.8,-77.4);
\draw (41.4,-91.3) -- (17,-95.5);
\draw (60.8,-105) -- (66.1,-119.3);
\end{tikzpicture}
\end{document}
